\section{Постановка задач}

В ходе лабораторной работы \textit{по варианту №7} для линейной регрессии вида 
\[
    f(x) = w_0 + w_1 x + w_2 x^2
\]

будут выполнены следующие задачи:

\begin{enumerate}
  \item Найти коэффициенты линейной регрессии \textit{методом наименьших квадратов (МНК)}.
  \item Вычислить среднеквадратическое отклонение (RMSE).
  \item Изобразить на графике точки выборки $(x, t)$ с графиком линейной регрессии.
\end{enumerate}

Для анализа предоставлена следующая выборка:

\begin{table}[h!]
\centering
\begin{tabular}{|c|c|c||c|c|c||c|c|c||c|c|c|}
\hline
№ & $x$ & $t$ & № & $x$ & $t$ & № & $x$ & $t$ & № & $x$ & $t$ \\ \hline \hline
1 & 1.0 & 2.0 & 6 & 1.5 & 1.3 & 11 & 2.0 & 1.0 & 16 & 2.5 & 1.3 \\ \hline
2 & 1.1 & 1.7 & 7 & 1.6 & 1.1 & 12 & 2.1 & 1.1 & 17 & 2.6 & 1.2 \\ \hline
3 & 1.2 & 1.8 & 8 & 1.7 & 1.2 & 13 & 2.2 & 1.0 & 18 & 2.7 & 1.5 \\ \hline
4 & 1.3 & 1.4 & 9 & 1.8 & 1.1 & 14 & 2.3 & 1.2 & 19 & 2.8 & 1.7 \\ \hline
5 & 1.4 & 1.4 & 10 & 1.9 & 1.0 & 15 & 2.4 & 1.1 & 20 & 2.9 & 1.9 \\ \hline
\end{tabular}
\caption{Исходные данные выборки}
\end{table}